\section{ANEXO F. ENTORNO DE INTEGRACIÓN CONTINUA (CI/CD) Y PUBLICACIÓN
AUTOMATIZADA}\label{anexo-f.-entorno-de-integraciuxf3n-continua-cicd-y-publicaciuxf3n-automatizada}

Para garantizar la integridad del software durante el ciclo de vida del
desarrollo y agilizar el flujo de despliegue, se ha diseñado e
implementado un pipeline de Integración Continua (CI) basado en
\textbf{GitHub Actions}. Esta estrategia de ingeniería de software
permite automatizar la compilación, verificación y empaquetado de la
aplicación en cada iteración, mitigando los riesgos asociados a la
integración manual de código y asegurando la disponibilidad inmediata de
artefactos listos para producción.

\subsection{ANEXO F.1. Arquitectura del Workflow y Disparadores
(Triggers)}\label{anexo-f.1.-arquitectura-del-workflow-y-disparadores-triggers}

El flujo de trabajo automatizado se define de manera declarativa
mediante la sintaxis YAML de GitHub Actions. Con el propósito de
optimizar los recursos de cómputo y mantener un control estricto sobre
la estabilidad de la rama principal, se han configurado dos disparadores
específicos:

\begin{itemize}
\item
  \textbf{Eventos de Empuje (\emph{Push}):} El pipeline se ejecuta de
  forma automática ante cualquier consolidación directa de código en la
  rama \texttt{main}, asegurando que cada incremento de software sea
  evaluado y empaquetado de inmediato.
\item
  \textbf{Solicitudes de Extracción (\emph{Pull Requests}):} La
  automatización actúa como una barrera de calidad (\emph{quality gate})
  ante cualquier intento de fusión hacia la rama \texttt{main}. Esto
  permite validar que las modificaciones propuestas por los
  desarrolladores no rompan el proceso de construcción del contenedor
  antes de que el código sea integrado definitivamente.
\end{itemize}

El entorno de ejecución seleccionado para los trabajos (\emph{jobs}) es
\texttt{ubuntu-latest}, lo que proporciona un entorno virtual limpio,
aislado y actualizado de manera nativa por la infraestructura de GitHub.

\subsection{ANEXO F.2. Desglose Técnico de las Etapas del
Pipeline}\label{anexo-f.2.-desglose-tuxe9cnico-de-las-etapas-del-pipeline}

El ciclo de vida del pipeline se divide en dos trabajos secuenciales y
dependientes, los cuales ejecutan tareas críticas de aprovisionamiento,
autenticación y despliegue:

\subsubsection{\texorpdfstring{ANEXO F.2.1. Trabajo de Construcción y
Publicación
(\texttt{create-docker-image})}{ANEXO F.2.1. Trabajo de Construcción y Publicación (create-docker-image)}}\label{anexo-f.2.1.-trabajo-de-construcciuxf3n-y-publicaciuxf3n-create-docker-image}

Es el núcleo técnico de la automatización y consta de los siguientes
pasos detallados:

\begin{enumerate}
\def\labelenumi{\arabic{enumi}.}
\item
  \textbf{Clonación del Repositorio y Gestión de Submódulos:} Se utiliza
  la acción oficial \texttt{actions/checkout@v2}. Debido a la
  arquitectura desacoplada del proyecto, es indispensable configurar el
  parámetro
  \texttt{submodules:\ \textquotesingle{}recursive\textquotesingle{}}.
  Esta directiva instruye al agente para que descargue e integre de
  manera automática el repositorio y código de \texttt{evmaudit} dentro
  de la estructura de directorios del pipeline.
\item
  \textbf{Autenticación en el Registro de Contenedores (GHCR):} Mediante
  la acción \texttt{docker/login-action@v2}, el pipeline establece una
  conexión segura con el registro oficial de GitHub (\texttt{ghcr.io}).
  El proceso se autentica de forma dinámica utilizando el actor del
  ciclo de vida (\texttt{\$\{\{\ github.actor\ \}\}}) y un token de
  acceso seguro almacenado de forma cifrada en los secretos del
  repositorio bajo la clave
  \texttt{\$\{\{\ secrets.IMAGES\_TFM\_UNIR\ \}\}}.
\item
  \textbf{Normalización del Espacio de Nombres y Doble Etiquetado
  (\emph{Multi-tagging}):} Las especificaciones de Docker y el estándar
  de la Open Container Initiative (OCI) prohíben estrictamente el uso de
  caracteres en mayúscula para los nombres de las imágenes de
  contenedores. Para solucionar esto, el pipeline ejecuta un script en
  Bash que convierte dinámicamente el nombre del repositorio a
  minúsculas utilizando el comando
  \texttt{tr\ \textquotesingle{}{[}:upper:{]}\textquotesingle{}\ \textquotesingle{}{[}:lower:{]}\textquotesingle{}}.
  Posteriormente, se implementa una estrategia de doble etiquetado para
  optimizar la trazabilidad y la inmutabilidad:

  \begin{itemize}
  \tightlist
  \item
    \textbf{Etiqueta por SHA (}\texttt{IMAGE\_SHA}\textbf{):} Vincula de
    forma unívoca el contenedor con el hash del \emph{commit} específico
    de Git que originó la compilación
    (\texttt{\$\{\{\ github.sha\ \}\}}). Esto permite realizar
    auditorías retrospectivas y despliegues deterministas en caso de
    fallos.
  \item
    \textbf{Etiqueta de Última Versión
    (}\texttt{IMAGE\_LATEST}\textbf{):} Sobrescribe el puntero
    \texttt{:latest} con la versión más reciente del software que haya
    superado la fase de construcción con éxito.
  \end{itemize}
\item
  \textbf{Construcción y \emph{Push} Multiplataforma:} Se invoca de
  manera directa el comando \texttt{docker\ build}, forzando la
  compilación bajo la arquitectura destino
  \texttt{-\/-platform\ linux/amd64} utilizando el \texttt{Dockerfile}
  del proyecto como plano de construcción. Una vez generadas las
  imágenes locales con sus respectivas etiquetas, se ejecutan las
  instrucciones de empuje (\emph{push}) hacia el registro seguro de
  GitHub, quedando el artefacto disponible para su consumo.
\end{enumerate}

\subsubsection{ANEXO F.2.2. Trabajo de Despliegue (deploy) y
Limitaciones del
Entorno}\label{anexo-f.2.2.-trabajo-de-despliegue-deploy-y-limitaciones-del-entorno}

Para garantizar una separación formal de conceptos en la arquitectura de
la integración, el pipeline implementa un segundo trabajo secuencial
denominado \texttt{deploy}. Utilizando la directiva
\texttt{needs:\ create-docker-image}, se establece una restricción de
dependencia estricta: esta etapa no puede inicializarse si el
empaquetado y la publicación previa de la imagen en el registro han
fallado.

En la fase actual del proyecto, \textbf{esta etapa automatizada no ha
sido implementada de forma activa y opera estrictamente como un punto de
anclaje (\emph{placeholder}) arquitectónico}. La justificación de esta
decisión de diseño radica en las limitaciones técnicas del entorno de
alojamiento seleccionado para las pruebas de concepto. La
infraestructura de EVMAudit se despliega externamente en la plataforma
PaaS \textbf{Railway} utilizando su modalidad de suscripción gratuita.
Este nivel de servicio impone restricciones en las interfaces de
programación (APIs) y en el uso de \emph{webhooks}, impidiendo la
ejecución de despliegues totalmente automatizados (\emph{Automated CD
Triggers}) desencadenados de forma directa mediante agentes de terceros
como GitHub Actions.

Por consiguiente, el flujo de trabajo en este punto se limita a
verificar la integridad de la secuencia de comandos en consola, quedando
el aprovisionamiento de la infraestructura supeditado a los mecanismos
nativos de la plataforma de destino. En el \textbf{siguiente apartado
(X.5. Despliegue de la Infraestructura)}, se expondrá de manera más
amplia la configuración, el aprovisionamiento y las características
operativas de dicho entorno en Railway.
